\uuid{ZCaa}
\titre{Dérivée sur un intervalle}
\auteur{bodin ; hanani ; mzari}
\organisation{exo7}
\licence{CC BY-NC-SA 4.0}
\datecreate{2020-01-20}
\niveau{L1}
\module{Analyse}
\chapitre{Dérivabilité des fonctions réelles}
\sousChapitre{Autre}
\difficulte{3}
\theme{Dérivées, dérivée sur un intervalle}
\exo7qcmid{lille-1.derivees.022}

\contenu{
\begin{question}
\qclassification{124.00}
\qtags{motcle=dérivée sur un intervalle}
Soit $f$ une fonction continue sur $[-1,1]$ telle que $f(0)=\pi$ et, pour tout $x\in ]-1,1[$, $\displaystyle f'(x)=\frac{1}{\sqrt{1-x^2}}$. Comment peut-on exprimer $f$ ?
\begin{answers}
    \bad{$f(x)=\sqrt{1-x^2}-1+\pi$}
    \good{$f(x)=\arcsin (x)+\pi$}
    \good{$\displaystyle f(x)=-\arccos x+\frac{3\pi}{2}$}
    \bad{Une telle fonction $f$ n'existe pas.}
\end{answers}
\begin{explanations}
On remarque que $f'(x)=(\arcsin x)'=(-\arccos x)'$. Donc, par continuité,
$$\forall x\in [-1,1],\; f(x)=\arcsin x+C_1=-\arccos x+C_2.$$
Mais $f(0)=\pi \Rightarrow C_1=\pi$ et $\displaystyle C_2=\frac{3\pi}{2}$.
\end{explanations}
\end{question}
}
